Для фиксированного графа предложенная модель решает прямую задачу: по
структуре работ и параметрам человеко-агентной ячейки оценивается
достижимое ускорение и строится расписание. Естественным продолжением
является обратная задача. Если в момент $t$ известна только часть графа
$G_t$, дальнейшую декомпозицию можно выбирать с учётом критического пути,
требуемого параллелизма и ограниченного ресурса человеческого внимания.
Тем самым модель становится не только средством оценки готового плана,
но и критерием формирования ещё не раскрытой части проекта.

Такую политику можно строить по принципу скользящего горизонта. После
результата, который обнаруживает новые работы или зависимости, обновляются
$G_t$, множество готовых задач и оценки фаз, после чего расписание незавершённой
части строится заново. Исследовательские задачи при этом можно представить
вершинами-разведками, результат которых раскрывает множество допустимых
продолжений $\mathcal D_t$. Для каждого $d\in\mathcal D_t$ можно использовать суррогатную оценку

\[
\Phi_t(d)=
\max\left\{
\frac{W_{4,t}(d)+\widehat Q_t(d)}{P},
L_{4,t}(d),
\gamma(P)H_t(d)
\right\},
\qquad
d_t^*\in\operatorname*{arg\,min}_{d\in\mathcal D_t}\Phi_t(d),
\]

где $\widehat Q_t(d)$ --- оценка будущего ожидания, не сводимого к уже известным
структурным границам. Вариант декомпозиции должен сохранять исходный объём
работ и критерий приёмки, а его дополнительные спецификации, проверки и интеграционные
издержки должны входить в $W_{4,t}(d)$ и $H_t(d)$. Если эмпирически известен
комфортный предел $P_{\max}^{\mathrm{attn}}$ одновременно сопровождаемых потоков, его превышение следует
запретить, а близость полезного параллелизма к $P_{\max}^{\mathrm{attn}}$ использовать как вторичный критерий
после минимизации $\Phi_t$. Построение допустимых преобразований графа, оценивание
$\widehat Q_t(d)$ и получение гарантий для такой адаптивной политики образуют отдельную задачу
совместного выбора декомпозиции и расписания; в настоящей работе этот критерий
приводится только как эскиз будущей теории.

Отдельное направление связано с распространением модели на другие архитектуры
многоагентных и роевых систем. Ближайшим к текущей постановке является
иерархический вариант, в котором агент-оркестратор заменяет разработчика во
внутреннем контуре: декомпозирует цель, ставит задачи рабочим агентам,
обслуживает их запросы и принимает результаты. Структурная аналогия здесь прямая:
оркестратор также является разделяемым ресурсом конечной мощности, а обращения
рабочих агентов образуют очередь и могут локально блокировать их потоки. При этом
человеческое внимание может сохраниться только на границах всего процесса --- при
постановке общей цели и итоговой приёмке.

Иная постановка возникает при алгоритмической оркестрации с аудитом и самопроверкой.
Макрограф задач может оставаться детерминированным, но каждая его вершина раскрывается
в локальный контур «исполнение --- аудит --- условная доработка --- повторный аудит». Если
такие контуры самодостаточны и не используют общий оркестратор, они не блокируют друг
друга, а внутренний человеческий ресурс мощности 1 исчезает. В терминах настоящей
модели аудит и доработка переходят в агентные фазы: $h_i$ для внутренних вершин
уменьшается, а $a_i$ растёт; знак изменения $k_i=a_i+h_i$ зависит от баланса
сэкономленной человеческой и добавленной агентной работы. При фиксированном числе проходов такой
контур можно развернуть в обычный фазовый DAG; условное завершение и неизвестное число
доработок потребуют стохастической модели повторной работы.

Более общее расширение должно описывать архитектуру не числом агентов, а графом
ресурсов и взаимодействий. Иерархические оркестраторы, децентрализованные рои,
взаимный аудит, специализированные планировщики и верификаторы, общая память и
гетерогенные ресурсы создают разные узкие места и разные формы координационных издержек.
Сравнение таких систем требует заменить звездообразную схему «один разработчик --- $P$
агентов» гетерогенной многоуровневой моделью расписаний и оценивать конфигурацию
целиком: топологию, правила маршрутизации, ресурсы аудита, качество результата и стоимость
дополнительных циклов самоконтроля. Несколько человеко-агентных ячеек требуют
дополнительных межъячеечных ресурсов и потому не сводятся к увеличению $P$ в M4.

Другим направлением является относительная и интервальная калибровка для организаций, не
располагающих историей точных измерений. В режиме холодного старта единицу $X$ можно трактовать
не как час, а как условную базовую задачу, после чего задавать $Z_i$, $k_i$ и $h_i$ экспертно
относительно неё. Если две конфигурации $A$ и $B$ описаны в одной шкале, а все их агентные и
человеческие фазы содержат один неизвестный общий множитель $\kappa>0$, то из
Предложения~\ref{prop:invariance} следует

\[
\frac{T_A^*(\kappa)}{T_B^*(\kappa)}=
\frac{T_A^*(1)}{T_B^*(1)}.
\]

Тем самым даже грубые количественные относительные оценки могут быть полезны
для выбора более быстрой схемы организации работ без обещания точного срока в часах.
Дальнейшее исследование должно заменить точечные оценки интервалами для $k_i$ и $h_i$,
определить условия устойчивого доминирования одного варианта над другим и выделить
случаи, в которых неопределённость относительного профиля способна изменить ранжирование.

Ещё одно направление --- стохастическая и робастная версия с повторной работой, отказами,
обучением и интервальными оценками $k$, $h$, $C$ и $\gamma$. Она должна давать не только
ожидаемый срок, но и квантили, вероятность нарушения дедлайна и условия устойчивого
доминирования одной конфигурации над другой.

Наконец, необходима эмпирическая калибровка на трассах реальной программной работы. Она
должна отдельно измерять автономные и человеческие интервалы, ожидание, повторную работу и
качество результата, чтобы проверить форму $C(P)$ и $\gamma(P)$ и отделить эффективный
параллелизм от номинального числа запущенных агентов.
